%=====================================================================
\chapter{Mathematics for Machine Learning}\label{ch:math}
%=====================================================================
\locator{1}{Part I --- Foundations}

\begin{outcomes}
By the end of this chapter you will be able to:
\begin{tight}
  \item \textbf{Compute} dot products, norms and distances, and write a linear
        model as a dot product.
  \item \textbf{Explain} what the product $X\mathbf{w}$ computes for a whole
        dataset at once.
  \item \textbf{Differentiate} a simple loss function and \textbf{perform}
        gradient descent steps by hand, including choosing a learning rate.
  \item \textbf{Apply} Bayes' theorem to a base-rate problem.
  \item \textbf{Estimate} a parameter by maximum likelihood, and \textbf{compute}
        the entropy of a distribution.
\end{tight}
\end{outcomes}

\begin{prereq}
School algebra, and the probability you used in \emph{Artificial Intelligence}.
\chref{introduction} for the idea of a model with weights $w_0, \dots, w_d$.
Appendix~B collects every formula in this chapter on one page.
\end{prereq}

\begin{keyterms}
vector $\bullet$ dot product $\bullet$ norm $\bullet$ Euclidean distance
$\bullet$ cosine similarity $\bullet$ matrix $\bullet$ transpose $\bullet$
derivative $\bullet$ partial derivative $\bullet$ gradient $\bullet$ chain rule
$\bullet$ gradient descent $\bullet$ learning rate $\bullet$ conditional
probability $\bullet$ Bayes' theorem $\bullet$ expectation $\bullet$ variance
$\bullet$ Gaussian $\bullet$ likelihood $\bullet$ maximum likelihood $\bullet$
entropy
\end{keyterms}

\section{Four Kinds of Mathematics, Four Jobs}

Machine learning uses four pieces of mathematics, and each does one job. It is
worth knowing which is which, because when a formula appears later you will know
what it is \emph{for} before you know what it says.

\rowcolors{2}{white}{lightbg}
\begin{longtable}{@{}L{3.0cm}L{5.4cm}L{6.6cm}@{}}
\toprule
\rowcolor{primary!15}
\textbf{Mathematics} & \textbf{Its job} & \textbf{Where it appears} \\
\midrule
\endhead
Linear algebra & \emph{Represent} examples and models & Every model: an example is
a vector, a dataset is a matrix, a linear model is a dot product (Chs 5, 6, 12, 13,
16) \\
Calculus & \emph{Optimise}: find the weights that minimise a loss & Gradient
descent in Chs 5, 6 and 13; backpropagation is the chain rule \\
Probability & \emph{Reason under uncertainty} & Naive Bayes (Ch 10), EM (Ch 17),
HMMs (Ch 18), every stochastic environment in Part V \\
Information theory & \emph{Measure surprise} & Decision-tree splits (Ch 9) and the
cross-entropy loss (Ch 6) \\
\bottomrule
\end{longtable}

\section{Vectors and the Dot Product}

\begin{definitionbox}[Vector, Dot Product and Norm]
A \term{vector} $\mathbf{x} = (x_1, \dots, x_d)$ is an ordered list of $d$
numbers --- in this course, usually the features of one example.\\[3pt]
The \term{dot product} of two vectors of the same length multiplies them
element by element and adds up:
\[
\mathbf{w} \cdot \mathbf{x} = \mathbf{w}^{\top}\mathbf{x} = \sum_{j=1}^{d} w_j x_j
\]
The \term{norm} (length) of a vector is $\|\mathbf{x}\| = \sqrt{\mathbf{x}\cdot\mathbf{x}}
= \sqrt{\sum_j x_j^2}$, and the \term{Euclidean distance} between two vectors is
the norm of their difference, $\|\mathbf{x} - \mathbf{z}\|$.
\end{definitionbox}

The dot product is the most important operation in this course. A linear model
--- linear regression, logistic regression, a perceptron, an SVM --- computes one
dot product and adds a bias: $\hat{y} = \mathbf{w}\cdot\mathbf{x} + b$. Each
weight says how much its feature pushes the prediction up or down.

\begin{worked}[One Task as a Vector]
A task has requirement clarity $0.9$ (the fraction of its acceptance criteria
written down), test coverage $0.6$ and $4$ open dependencies:
$\mathbf{x} = (0.9,\ 0.6,\ 4)$. A model has weights $\mathbf{w} = (2,\ 3,\ 0.25)$
and bias $b = -3$.

\smallskip
\textbf{Score:} $\mathbf{w}\cdot\mathbf{x} + b = (2)(0.9) + (3)(0.6) + (0.25)(4) - 3
= 1.8 + 1.8 + 1.0 - 3 = \mathbf{1.6}$.

\smallskip
\textbf{Distance to a second task} $\mathbf{z} = (0.5,\ 0.4,\ 1)$:
\[
\|\mathbf{x}-\mathbf{z}\| = \sqrt{0.4^2 + 0.2^2 + 3^2} = \sqrt{0.16 + 0.04 + 9}
= \sqrt{9.2} \approx \mathbf{3.03}
\]
\textbf{Look at where the distance came from.} Of the $9.2$ under the root, $9$
is dependencies. Clarity and coverage --- arguably the more important differences
--- contribute almost nothing, simply because they are measured on a scale of 0
to 1 and dependencies are not. Any method built on distance (\chref{knn},
\chref{clustering}) needs the features put on a common scale first
(\chref{workflow}).
\end{worked}

\begin{conceptbox}[Cosine Similarity: Direction Without Length]
Sometimes length is a nuisance. Two tickets reporting the same fault, one ten
times longer, have word-count vectors pointing the same way but of very different
lengths. \term{Cosine similarity} compares directions only:
$\cos\theta = \dfrac{\mathbf{x}\cdot\mathbf{z}}{\|\mathbf{x}\|\,\|\mathbf{z}\|}$,
which is $1$ for vectors pointing the same way, $0$ for perpendicular ones and
$-1$ for opposite ones.
\end{conceptbox}

\section{Matrices: A Whole Dataset at Once}

\begin{definitionbox}[Data Matrix and Matrix--Vector Product]
Stack $n$ examples as the rows of a \term{matrix} $X$ with $n$ rows and $d$
columns. Then $X\mathbf{w}$ is a vector of $n$ numbers: row $i$ of the result is
the dot product of example $i$ with $\mathbf{w}$. The \term{transpose}
$X^{\top}$ swaps rows and columns.
\end{definitionbox}

\dsfig{%
\begin{tikzpicture}[font=\scriptsize,
  cell/.style={draw=primary!50, fill=primary!6, minimum width=0.9cm, minimum height=0.5cm,
               font=\scriptsize},
  wcell/.style={draw=goldhl, fill=goldhl!10, minimum width=0.9cm, minimum height=0.5cm,
                font=\scriptsize},
  ycell/.style={draw=greenhl, fill=greenhl!10, minimum width=1.2cm, minimum height=0.5cm,
                font=\scriptsize}]
% X
\foreach \r/\a/\b/\c in {0/0.9/0.6/4, 1/0.5/0.4/1, 2/0.7/0.9/6}{
  \node[cell] at (0,-\r*0.5) {\a};
  \node[cell] at (0.9,-\r*0.5) {\b};
  \node[cell] at (1.8,-\r*0.5) {\c};}
\fill[accent!15] (-0.45,0.25) rectangle (2.25,-0.25);
\foreach \a/\x in {0.9/0,0.6/0.9,4/1.8}{\node[cell, fill=accent!12, draw=accent] at (\x,0) {\a};}
\node[font=\scriptsize\bfseries, text=primary] at (0.9,0.65) {$X$ \quad ($3 \times 3$)};
\node[font=\tiny, text=gray!55!black, anchor=east] at (-0.55,0) {task 1};
\node[font=\tiny, text=gray!55!black, anchor=east] at (-0.55,-0.5) {task 2};
\node[font=\tiny, text=gray!55!black, anchor=east] at (-0.55,-1.0) {task 3};
\node[font=\large] at (2.75,-0.5) {$\times$};
% w
\foreach \r/\a in {0/2, 1/3, 2/0.25}{\node[wcell] at (3.5,-\r*0.5) {\a};}
\node[font=\scriptsize\bfseries, text=goldhl!50!black] at (3.5,0.65) {$\mathbf{w}$};
\node[font=\large] at (4.3,-0.5) {$+\,b$};
\node[font=\large] at (5.1,-0.5) {$=$};
% y
\foreach \r/\a in {0/1.60, 1/$-0.55$, 2/2.60}{\node[ycell] at (6.0,-\r*0.5) {\a};}
\node[ycell, draw=accent, fill=accent!12] at (6.0,0) {1.60};
\node[font=\scriptsize\bfseries, text=greenhl!50!black] at (6.0,0.65) {$\hat{\mathbf{y}}$};
% note
\node[anchor=north west, align=left, text width=5.4cm, font=\scriptsize, text=gray!55!black]
     at (7.0,0.55)
     {\textbf{One multiplication, every prediction.} The highlighted row of $X$ dotted
      with $\mathbf{w}$, plus $b=-3$, gives the highlighted entry of $\hat{\mathbf{y}}$ ---
      the score computed by hand in the worked example. Libraries such as NumPy
      compute all $n$ rows at once, which is why models are written with matrices.};
\end{tikzpicture}}%
{A matrix--vector product scores every example in one operation. Each row of $X$ is
one example; each entry of $\hat{\mathbf{y}}$ is that row's dot product with the
weights.}{matvec}

You will meet one more matrix operation, in \chref{linreg}: the \term{inverse}
$A^{-1}$ of a square matrix, which undoes multiplication by $A$ the way dividing by
3 undoes multiplying by 3. You will never compute one by hand; NumPy does.

\section{Derivatives and Gradients}

\begin{definitionbox}[Derivative, Partial Derivative, Gradient]
The \term{derivative} $\dfrac{dL}{dw}$ of a function $L(w)$ is its slope: how much
$L$ changes per unit increase in $w$. For a function of several weights, the
\term{partial derivative} $\dfrac{\partial L}{\partial w_j}$ is the slope in the
direction of $w_j$ alone, holding the others fixed. The \term{gradient} collects
them into a vector:
\[
\nabla L = \left(\frac{\partial L}{\partial w_1}, \dots, \frac{\partial L}{\partial w_d}\right)
\]
The gradient points in the direction in which $L$ increases fastest.
\end{definitionbox}

You need only four rules, and Appendix~B lists them: the derivative of $w^k$ is
$k\,w^{k-1}$; of a constant is 0; of a sum is the sum of the derivatives; and the
\term{chain rule}: if $L = f(g(w))$ then $\dfrac{dL}{dw} = f'(g(w))\cdot g'(w)$.
For example, $L(w) = (w-3)^2$ has outer function $u^2$ and inner function
$u = w - 3$, so $\dfrac{dL}{dw} = 2(w-3)\cdot 1 = 2(w-3)$. The chain rule is the
whole of backpropagation (\chref{ann}).

\section{Gradient Descent}

Training a model means finding weights that make a \term{loss} $L(\mathbf{w})$
small. For most models there is no formula for the best weights, so we walk
downhill instead.

\begin{definitionbox}[Gradient Descent]
Start from any weights. Repeatedly take a small step against the gradient:
\[
\mathbf{w} \leftarrow \mathbf{w} - \alpha\, \nabla L(\mathbf{w})
\]
The \term{learning rate} $\alpha > 0$ sets the step size. Stop when the loss stops
decreasing.
\end{definitionbox}

\dsfig{%
\begin{tikzpicture}
\begin{groupplot}[
  group style={group size=2 by 1, horizontal sep=1.4cm},
  width=6.8cm, height=5.0cm,
  axis lines=left, tick label style={font=\tiny}, label style={font=\tiny},
  title style={font=\scriptsize\bfseries}, xlabel={$w$}, ylabel={$L(w)=(w-3)^2$},
  xmin=-2.5, xmax=9, ymin=0, ymax=36]
\nextgroupplot[title={$\alpha=0.1$: steady descent}]
\addplot[gray!60, line width=1pt, domain=-2.5:9, samples=80] {(x-3)^2};
\addplot[greenhl!70!black, mark=*, mark size=1.8pt, line width=0.9pt]
  coordinates {(0,9)(0.6,5.76)(1.08,3.6864)(1.464,2.3593)(1.7712,1.5099)};
\node[font=\tiny, greenhl!45!black, anchor=south west] at (axis cs:0.1,9.3) {start $w=0$};
\draw[accent, dashed] (axis cs:3,0) -- (axis cs:3,6);
\node[font=\tiny, accent, anchor=south] at (axis cs:3,6) {minimum};
\nextgroupplot[title={$\alpha=1.1$: each step overshoots further}]
\addplot[gray!60, line width=1pt, domain=-2.5:9, samples=80] {(x-3)^2};
\addplot[accent, mark=*, mark size=1.8pt, line width=0.9pt]
  coordinates {(0,9)(6.6,12.96)(-1.32,18.6624)(8.184,26.874)};
\node[font=\tiny, accent, anchor=south west] at (axis cs:0.1,9.3) {start};
\draw[accent, dashed] (axis cs:3,0) -- (axis cs:3,6);
\end{groupplot}
\end{tikzpicture}}%
{Gradient descent on $L(w)=(w-3)^2$ from $w=0$. With a small learning rate each step
closes a fixed fraction of the gap; with one that is too large, each step jumps past the
minimum by more than the last, and the weights diverge.}{gradient-descent}

\Needspace{16\baselineskip}
\begin{worked}[Four Steps of Gradient Descent]
Minimise $L(w) = (w-3)^2$ starting from $w = 0$ with $\alpha = 0.1$. The
gradient is $L'(w) = 2(w-3)$.

\smallskip
\begin{center}
\begin{tabular}{@{}cccc@{}}
\toprule
\textbf{Step} & $w$ & $L'(w) = 2(w-3)$ & $w - 0.1\,L'(w)$ \\
\midrule
0 & 0 & $-6$ & $0 + 0.6 = 0.6$ \\
1 & 0.6 & $-4.8$ & $0.6 + 0.48 = 1.08$ \\
2 & 1.08 & $-3.84$ & $1.08 + 0.384 = 1.464$ \\
3 & 1.464 & $-3.072$ & $1.464 + 0.3072 = 1.7712$ \\
\bottomrule
\end{tabular}
\end{center}

\smallskip
\textbf{The pattern.} The gap to the minimum goes $3 \to 2.4 \to 1.92 \to 1.536
\to 1.2288$: every step multiplies it by $1 - 2\alpha = 0.8$. So
\begin{tight}
  \item $\alpha = 0.1$: the gap shrinks by 20\% a step --- steady.
  \item $\alpha = 0.5$: the factor is 0, so one step lands exactly on 3.
  \item $\alpha = 1.1$: the factor is $-1.2$ --- the weight flips to the far side
        and lands further away each time: $0 \to 6.6 \to -1.32 \to 8.18$. It
        diverges (\dsref{gradient-descent}, right).
\end{tight}
Real losses are not this simple, and the right $\alpha$ is found by trying
several (\chref{generalization}), but the three behaviours --- slow, just right,
divergent --- are exactly the ones you will see.
\end{worked}

\section{Probability}

\begin{definitionbox}[Conditional Probability and Bayes' Theorem]
The \term{conditional probability} of $A$ given $B$ is
$P(A \mid B) = \dfrac{P(A \text{ and } B)}{P(B)}$. Reversing the condition gives
\term{Bayes' theorem}:
\[
P(A \mid B) = \frac{P(B \mid A)\, P(A)}{P(B)},
\qquad P(B) = P(B\mid A)P(A) + P(B \mid \text{not }A)\,P(\text{not }A)
\]
$P(A)$ is the \term{prior}, $P(B\mid A)$ the \term{likelihood} and $P(A \mid B)$ the
\term{posterior}. Two events are \term{independent} if $P(A \text{ and } B) =
P(A)\,P(B)$.
\end{definitionbox}

\begin{worked}[A Blocked Release You Should Not Panic About]
An automated release check catches 99\% of builds that would break production,
and wrongly blocks 1\% of good builds. One build in a thousand is broken. The
check blocks a build. What is the probability that it really is broken?

\smallskip
Let $A$ = broken build, $B$ = blocked. $P(A) = 0.001$, $P(B\mid A) = 0.99$,
$P(B \mid \text{not }A) = 0.01$.
\[
P(B) = (0.99)(0.001) + (0.01)(0.999) = 0.00099 + 0.00999 = 0.01098
\]
\[
P(A \mid B) = \frac{0.00099}{0.01098} \approx \mathbf{0.090}
\]
\textbf{Only 9\% of blocked builds are broken}, from a check that is ``99\%
accurate''. There are so many more good builds that their 1\% of false blocks
outnumbers the real failures ten to one. This is the \term{base-rate}
effect, and it is why \chref{evaluation} insists on precision rather than
accuracy. Naive Bayes (\chref{bayes}) is this calculation, done for every class.
\end{worked}

\section{Expectation, Variance and the Gaussian}

The \term{expectation} (mean) of a random variable is its probability-weighted
average, $\mathbb{E}[X] = \sum_x x\,P(x)$; from a sample it is estimated by the
sample mean $\bar{x}$. The \term{variance} is the expected squared distance from
the mean, $\operatorname{Var}(X) = \mathbb{E}[(X - \mathbb{E}[X])^2]$, and its
square root is the standard deviation $\sigma$. For the values $2, 4, 4, 4, 5, 5,
7, 9$ the mean is 5, the squared deviations sum to 32, the variance is
$32/8 = 4$ and $\sigma = 2$.

The \term{Gaussian} (normal) distribution with mean $\mu$ and variance
$\sigma^2$ has density
\[
p(x) = \frac{1}{\sqrt{2\pi\sigma^2}} \exp\!\left(-\frac{(x-\mu)^2}{2\sigma^2}\right)
\]
It appears in Gaussian naive Bayes (\chref{bayes}) and Gaussian mixtures
(\chref{semisup}), and whenever a model assumes that its errors are small,
symmetric and unsurprising.

\section{Maximum Likelihood}

Most learning algorithms can be read as answering one question: \emph{which
parameters make the observed data most probable?}

\begin{definitionbox}[Likelihood and Maximum Likelihood Estimation]
For data $D$ and parameters $\theta$, the \term{likelihood} $\mathcal{L}(\theta) =
P(D \mid \theta)$ is the probability of the data under those parameters, read as a
function of $\theta$. The \term{maximum likelihood estimate}
$\hat{\theta}_{ML}$ is the $\theta$ that maximises it. In practice one maximises
the \term{log-likelihood} $\log \mathcal{L}(\theta)$ instead: the logarithm turns
products into sums and has the same maximum.
\end{definitionbox}

\dsfig{%
\begin{tikzpicture}
\begin{axis}[width=8.4cm, height=4.6cm, axis lines=left,
  xlabel={$p$, the probability that a sprint delivers}, ylabel={likelihood $p^7(1-p)^3$},
  tick label style={font=\tiny}, label style={font=\tiny},
  xmin=0, xmax=1, ymin=0, ymax=0.0025, scaled y ticks=false,
  yticklabel style={/pgf/number format/fixed, /pgf/number format/precision=4}]
\addplot[primary, line width=1.3pt, domain=0:1, samples=120] {x^7*(1-x)^3};
\draw[accent, dashed, line width=0.9pt] (axis cs:0.7,0) -- (axis cs:0.7,0.002224);
\node[font=\tiny, accent, anchor=west] at (axis cs:0.72,0.00228) {$\hat{p}_{ML}=0.7$};
\node[font=\tiny, gray!55!black, anchor=west, align=left] at (axis cs:0.03,0.0019)
     {7 delivered, 3 missed:\\$p=0.5$ explains the data\\less well than $p=0.7$};
\end{axis}
\end{tikzpicture}}%
{The likelihood of seven delivered sprints in ten, as a function of $p$. It is largest
at $p=0.7$ --- the observed proportion, which is what maximum likelihood gives for any
such count.}{likelihood}

\begin{worked}[How Often Does the Team Deliver?]
A team delivered everything it committed to in 7 of its last 10 sprints. Let $p$
be the probability that a sprint delivers, and treat sprints as independent. The
likelihood of that record is $\mathcal{L}(p) = p^7(1-p)^3$, and the log-likelihood is
$\ell(p) = 7\ln p + 3\ln(1-p)$. Set its derivative to zero:
\[
\frac{d\ell}{dp} = \frac{7}{p} - \frac{3}{1-p} = 0
\;\Rightarrow\; 7(1-p) = 3p \;\Rightarrow\; \hat{p} = \mathbf{0.7}
\]
The answer is what intuition says --- the observed fraction --- but the method
is general: the same argument gives the least-squares line of \chref{linreg} and
the cross-entropy loss of \chref{logreg}, and the EM algorithm of
\chref{semisup} is a way of doing it when some data is missing.
\end{worked}

\begin{alertbox}[Maximum Likelihood Trusts Small Samples Too Much]
One sprint, one success: $\hat{p}_{ML} = 1$, a team that can never miss. A word
that never appeared in the training tickets gets probability zero, and a single
zero makes a whole product zero. \chref{bayes} fixes this with
\term{Laplace smoothing}: add one imaginary observation of every outcome.
\end{alertbox}

\section{Information and Entropy}

\begin{definitionbox}[Entropy]
The \term{entropy} of a distribution $(p_1, \dots, p_k)$ measures its uncertainty,
in bits:
\[
H = -\sum_{i=1}^{k} p_i \log_2 p_i \qquad (\text{with } 0 \log_2 0 = 0)
\]
It is 0 when one outcome is certain, and largest --- $\log_2 k$ bits --- when all
$k$ outcomes are equally likely.
\end{definitionbox}

\begin{worked}[Entropy of a Set of Examples]
Fourteen past sprints, of which 9 met their sprint goal and 5 did not:
\[
H = -\tfrac{9}{14}\log_2\tfrac{9}{14} - \tfrac{5}{14}\log_2\tfrac{5}{14}
  = 0.410 + 0.531 = \mathbf{0.940} \text{ bits}
\]
A 7--7 split would give exactly 1 bit; a 14--0 split, 0 bits. A decision tree
(\chref{trees}) chooses the question that lowers this number the most --- that
is, the question whose answer leaves the least uncertainty behind.
\end{worked}

\begin{pitfall}
\begin{tight}
  \item \textbf{Forgetting the minus sign in gradient descent.} $\mathbf{w} +
        \alpha\nabla L$ climbs the loss. The update subtracts.
  \item \textbf{Reading a diverging loss as ``needs more steps''.} If the loss
        grows, the learning rate is too large; more steps make it worse.
  \item \textbf{Confusing $P(A\mid B)$ with $P(B \mid A)$.} ``The check
        catches 99\% of broken builds'' is $P(\text{blocked}\mid\text{broken})$;
        the release manager wants $P(\text{broken}\mid\text{blocked})$, which was
        9\%.
  \item \textbf{Computing distances on unscaled features.} The feature with the
        largest units decides the answer.
  \item \textbf{Multiplying many small probabilities.} The product underflows
        to zero on a computer; add log-probabilities instead.
\end{tight}
\end{pitfall}

%---------------------------------------------------------------------
\begin{lab}[The Chapter in NumPy]
\begin{lstlisting}[language=Python]
import numpy as np

# --- vectors and the dot product -----------------------------------------
x = np.array([0.9, 0.6, 4.0])
w = np.array([2.0, 3.0, 0.25])
print("score w.x + b :", round(w @ x - 3, 2))            # 1.6
z = np.array([0.5, 0.4, 1.0])
print("distance      :", round(np.linalg.norm(x - z), 2))  # 3.03

# --- a whole dataset at once ---------------------------------------------
X = np.array([[0.9, 0.6, 4], [0.5, 0.4, 1], [0.7, 0.9, 6]])
print("X w + b       :", X @ w - 3)                      # [1.6 -0.55 2.6]

# --- gradient descent on L(w) = (w - 3)^2 --------------------------------
for alpha in (0.1, 0.5, 1.1):
    wt = 0.0
    path = [wt]
    for step in range(4):
        grad = 2 * (wt - 3)
        wt = wt - alpha * grad
        path.append(round(wt, 3))
    print(f"alpha={alpha}: {path}")

# --- Bayes by simulation: how many blocked builds are broken? ------------
rng = np.random.default_rng(0)
n = 1_000_000
broken = rng.random(n) < 0.001
blocked = np.where(broken, rng.random(n) < 0.99, rng.random(n) < 0.01)
print("P(broken | blocked), simulated:",
      round(broken[blocked].mean(), 3))                      # about 0.09

# --- entropy -------------------------------------------------------------
def entropy(counts):
    p = np.array(counts) / np.sum(counts)
    p = p[p > 0]
    return -(p * np.log2(p)).sum()
print("H(9 yes, 5 no) :", round(entropy([9, 5]), 3))       # 0.94
print("H(7 yes, 7 no) :", round(entropy([7, 7]), 3))       # 1.0
\end{lstlisting}
Change \texttt{0.001} in the simulation to \texttt{0.05} --- a codebase where
broken builds are common --- and predict the new answer before you run it.
\end{lab}

%---------------------------------------------------------------------
\begin{checkpoint}
\begin{tightnum}
  \item Compute the Euclidean distance between $(1, 2)$ and $(4, 6)$.
  \item Take one gradient descent step on $L(w) = (w+1)^2$ from $w = 2$ with
        $\alpha = 0.25$. Is the new $w$ closer to the minimum?
  \item Compute the entropy of the distribution $(0.5,\ 0.25,\ 0.25)$.
\end{tightnum}
\end{checkpoint}

%---------------------------------------------------------------------
\begin{chaptersummary}
\begin{tight}
  \item An example is a vector, a dataset is a matrix, and a linear model is a
        dot product plus a bias; $X\mathbf{w}$ scores every example at once.
  \item Distances are dominated by the features with the largest units, so scale
        features before comparing examples.
  \item The gradient points uphill; gradient descent steps against it. Too small
        a learning rate is slow, too large diverges.
  \item Bayes' theorem reverses a conditional probability. When the prior is
        small, even an accurate test produces mostly false alarms.
  \item Maximum likelihood chooses the parameters that make the data most
        probable; work with the log-likelihood.
  \item Entropy measures uncertainty in bits: zero when the outcome is certain,
        largest when every outcome is equally likely.
\end{tight}
\end{chaptersummary}

\begin{reviewq}
  \item \textbf{(MCQ)} The dot product of $(1, -1, 2)$ and $(3, 1, 1)$ is:
        (a)~2 (b)~4 (c)~6 (d)~0.
  \item \textbf{(MCQ)} If the loss increases after every gradient descent step,
        the most likely cause is: (a)~too few steps (b)~a learning rate that is
        too large (c)~too much data (d)~a learning rate that is too small.
  \item \textbf{(MCQ)} A new ticket is equally likely to belong to any of six
        categories. The entropy of its category is: (a)~1 bit
        (b)~$\log_2 6 \approx 2.58$ bits (c)~6 bits (d)~0 bits.
  \item \textbf{(Short)} Explain in words what the gradient of a loss function
        tells you, and why gradient descent subtracts it.
  \item \textbf{(Short)} Why do we maximise the log-likelihood rather than the
        likelihood?
  \item \textbf{(Short)} Give an example, not from this chapter, where
        $P(A\mid B)$ and $P(B \mid A)$ are very different.
  \item \textbf{(Applied)} A code-review bot flags 2\% of clean pull requests
        and 95\% of those that contain a critical defect; 0.2\% of pull requests
        contain one. Compute the probability that a flagged pull request contains
        a critical defect, and state what this means for the reviewers.
  \item \textbf{(Applied)} Perform three gradient descent steps on
        $L(w) = 2(w-5)^2$ from $w = 1$ with $\alpha = 0.1$. Find the factor by
        which the gap to the minimum shrinks each step, and the largest
        learning rate for which descent still converges.
\end{reviewq}
